% Generated draft: requires mathematical and source review.
\documentclass[11pt]{article}
\usepackage[margin=25mm]{geometry}
\usepackage{fontspec}
\setmainfont{Latin Modern Roman}
\usepackage{amsmath,amssymb,mathtools,unicode-math}
\setmathfont{Latin Modern Math}
\usepackage{enumitem,needspace,xurl,hyperref}
\setlength{\parindent}{0pt}
\setlength{\parskip}{5pt}
\setlength{\emergencystretch}{3em}
\newcommand{\N}{\mathbb{N}}
\newcommand{\Z}{\mathbb{Z}}
\newcommand{\Q}{\mathbb{Q}}
\newcommand{\R}{\mathbb{R}}
\newcommand{\C}{\mathbb{C}}
\newcommand{\F}{\mathbb{F}}
\newcommand{\A}{\mathbb{A}}
\newcommand{\PP}{\mathbb{P}}
\newcommand{\E}{\mathbb{E}}
\newcommand{\eps}{\varepsilon}
\newcommand{\dd}{\,\mathrm{d}}
\begin{document}
\section*{UnsolvedMath ID: 30002162}
\textbf{Source problem code:} catalogue entry
\subsection{Definitions and mathematical statement}
Put \(s=(d+2)/6\), where \(d\) is the spatial dimension referred to in the source; the intended range of \(d\) is not specified. For a vector field \(v\colon\R^d\to\R^d\), use the standard inhomogeneous Sobolev space
\[
H^s(\R^d;\R^d)=\left\{v:\sum_{j=1}^d\int_{\R^d}(1+|\xi|^2)^s|\widehat v_j(\xi)|^2\,\dd\xi<\infty\right\},
\]
with the displayed integral sum defining the squared norm. Here \(\widehat v_j\) denotes the Fourier transform of the \(j\)-th component. This makes explicit the conventional interpretation of the source notation \(H^s(\R^d)\).

\textbf{Source question, with unresolved hypotheses.} Let \(u\) be a solution of the intended \(d\)-dimensional hyperdissipative flow on \(\R^d\), regular before a possible finite singular time \(t_*\). Must every such finite-time blowup satisfy
\[
\lim_{t\uparrow t_*}\|u(t)\|_{H^{(d+2)/6}(\R^d)}^2=+\infty?
\]
The supplied background identifies the intended setting as hyperdissipative Navier--Stokes and calls this an endpoint lower-branch regularity criterion. However, the record does not specify the evolution equation, dissipation operator or exponent, admissible dimensions, initial-data assumptions, solution class, or meaning of regularity and blowup. Accordingly, the quantifier over solutions remains conditional on these missing hypotheses.

The source also asks whether finiteness of this norm on \([0,T]\) suffices for regularity, and describes this as equivalent. A possible precise interpretation is that, for every finite \(T>0\), a solution regular on \([0,T)\) satisfying
\[
\sup_{0\le t<T}\|u(t)\|_{H^{(d+2)/6}(\R^d)}<\infty
\]
remains regular through \(T\). This is an interpretation requiring review: the source does not distinguish uniform boundedness from pointwise finiteness, or specify whether a value at \(T\) is available. Moreover, this boundedness criterion by itself excludes a finite bound near blowup; identifying it with the stated full limit to infinity requires additional solution-theoretic justification absent from the record.
\subsection{Short English statement}
The question asks whether every finite-time singularity of the intended hyperdissipative flow forces the \(H^{(d+2)/6}(\R^d)\) norm to tend to infinity. The equation, dissipation parameters, solution hypotheses, and interpretation of the proposed equivalent regularity criterion need clarification.
\subsection{Sources}
\begin{itemize}
\item Local UnsolvedMath registry, supplied complete source record: permanent ID \(30002162\), exact source problem code \texttt{catalogue entry}, title \emph{Borderline Sobolev Blowup Criteria for Hyperdissipative Flow}; collection: Oberwolfach Reports Open Problems. The project identifies \texttt{data/unsolvedmath/problems.json} and \texttt{catalog/problems.json} as registry sources; the supplied record does not identify which contains this entry.
\item Supplied report citation: \emph{Mathematical Aspects of Hydrodynamics (2013)}, Oberwolfach Reports, pp.\ 2389--2427, DOI \texttt{10.4171/owr/2012/39}, \url{https://doi.org/10.4171/owr/2012/39}. The literature-source label is ``OWR report 39.'' The record assigns proposed year 2012 but names no individual proposer. The relationship between that year, the DOI year, and the citation's 2013 date is not explained.
\item Supplied literature reference labeled ``Escauriaza--Seregin--Sverak'': \url{https://doi.org/10.1070/RM2003v058n02ABEH000610}. No title or explanation of its precise role in this formulation is supplied.
\item Literature assessment dated exactly \texttt{2026-08-21}: the supplied curation labels the problem open, describes the endpoint lower-branch criterion as unproved, and reports known stronger-dissipation regimes and special cases. This is supplied dated curation, not an independently verified literature assessment; the background requests independent verification of time-sensitive status and cited claims.
\item The registry supplies the label \texttt{corrected\_verified} and reports that the mathematical content was checked against the official Oberwolfach report, with notation, context, and question form normalized without changing the source-level claim. This records the supplied verification assertion; no source was independently visited for this draft.
\end{itemize}
\end{document}
